\textbf{Quiz-1}
\subsection{Variation Question 1.1}
Which of the following 8-bit binary words could represent $-22$ using the representations for negative numbers we saw in the lectures?

Tick the two correct answers.
\textbf{Correct Answer(s):}
\begin{itemize}
    \item 10010110
    \item 11101010
\end{itemize}

\textbf{Incorrect Answers:}
\begin{itemize}
    \item 00010110
    \item 11101001
    \item 11101000
\end{itemize}

\textbf{Solution \& Feedback:}\\
22 is 00010110 in 8-bit binary.\\
Therefore, in sign-and-magnitude representation $-22$ is 10010110.\\
For two’s complement, we flip all bits to get 11101001 and then add 1 to the LSB to get the two’s complement representation of $-22$: 11101010.

\subsection{Variation Question 1.2}

Which of the following 8-bit binary words could represent $-54$ using the representations for negative numbers we saw in the lectures?
Tick the two correct answers.
\textbf{Correct Answer(s):}
\begin{itemize}
    \item 10110110
    \item 11001010
\end{itemize}

\textbf{Incorrect Answers:}
\begin{itemize}
    \item 00110110
    \item 11001001
    \item 11001000
\end{itemize}

\textbf{Solution \& Feedback:}\\
54 is 00110110 in 8-bit binary.\\
Therefore, in sign-and-magnitude representation $-54$ is 10110110.\\
For two’s complement, we flip all bits to get 11001001 and then add 1 to the LSB to get the two’s complement representation of $-54$: 11001010.

\subsection{Variation Question 1.3}
Which of the following 8-bit binary words could represent $-38$ using the representations for negative numbers we saw in the lectures?
Tick the two correct answers.

\textbf{Correct Answer(s):}
\begin{itemize}
    \item 10100110
    \item 11011010
\end{itemize}
\textbf{Incorrect Answers:}
\begin{itemize}
    \item 00100110
    \item 11011001
    \item 11011000
\end{itemize}
\textbf{Solution \& Feedback:}\\
38 is 00100110 in 8-bit binary.\\
Therefore, in sign-and-magnitude representation $-38$ is 10100110.\\
For two’s complement, we flip all bits to get 11011001 and then add 1 to the LSB to get the two’s complement representation of $-38$: 11011010.

\subsection{Variation Question 2.1}
Which of the following statements correctly describe components or actions involved in the Fetch/Execute cycle?
Tick the two correct answers.\\
\textbf{Correct Answer(s):}
\begin{itemize}
    \item The Program Counter (PC) points to the next instruction to be fetched from memory.
    \item The ALU is responsible for performing computations on data.
\end{itemize}

\textbf{Incorrect Answers:}
\begin{itemize}
    \item The Control Unit performs arithmetic and logical operations.
    \item During the Data Fetch (DF) step, data is destroyed after being fetched from memory.
    \item The Instruction Fetch (IF) step is where the ALU executes an instruction.
\end{itemize}

\textbf{Solution \& Feedback:}\\
The Control Unit performs arithmetic and logical operations.\\
This is incorrect. It is ALU that performs such operations instead of the Control Unit.\\

During the Data Fetch (DF) step, data is destroyed after being fetched from memory.\\
This is incorrect. Data remains in the memory after the DF step.\\

The Instruction Fetch (IF) step is where the ALU executes an instruction.\\
This is incorrect. It should be Instruction Execution instead of Instruction Fetch.

\subsection{Variation Question 2.2}

Which of the following statements correctly describe components or actions involved in the Fetch/Execute cycle?

Tick the two correct answers.

\textbf{Correct Answer(s):}
\begin{itemize}
    \item During the Data Fetch (DF) step, data values remain unchanged in the memory.
    \item During the Result Return (RR) step, the result is stored back into memory.
\end{itemize}

\textbf{Incorrect Answers:}
\begin{itemize}
    \item The Program Counter (PC) points to the current instruction to be fetched from memory.
    \item The ISA performs arithmetic and logical operations.
    \item The Instruction Fetch (IF) step is where the ALU executes an instruction.
\end{itemize}

\textbf{Solution \& Feedback:}\\
The Program Counter (PC) points to the current instruction to be fetched from memory.\\
This is incorrect. The PC points to the next instruction instead of the current instruction.\\

The ISA performs arithmetic and logical operations.\\
This is incorrect. It is ALU that performs such operations.\\

The Instruction Fetch (IF) step is where the ALU executes an instruction.\\
This is incorrect. It should be Instruction Execution instead of Instruction Fetch for the step where the ALU executes an instruction.

\subsection{Variation Question 2.3}

Which of the following statements correctly describe components or actions involved in the Fetch/Execute cycle?

Tick the two correct answers.

\textbf{Correct Answer(s):}
\begin{itemize}
    \item The Instruction Decoder identifies the type of operation, and the operands involved in the instruction.
    \item During the Data Fetch (DF) step, data values remain unchanged in the memory.
\end{itemize}

\textbf{Incorrect Answers:}
\begin{itemize}
    \item The Program Counter (PC) points to the next instruction to be fetched from the register.
    \item After the Result Return (RR) step, the result is stored back into memory.
    \item The Instruction Fetch (IF) step is where the ALU executes an instruction.
\end{itemize}

\textbf{Solution \& Feedback:}\\
The Program Counter (PC) points to the next instruction to be fetched from the register.\\
This is incorrect. It is fetched from the memory instead of the register.\\

After the Result Return (RR) step, the result is stored back into memory.\\
This is incorrect. It should be ``During the RR step'' instead of ``After.''\\

The Instruction Fetch (IF) step is where the ALU executes an instruction.\\
This is incorrect. It should be Instruction Execution instead of Instruction Fetch.

\subsection{Variation Question 3.1}

Which of the following statements correctly describe aspects of the MIPS architecture and assembly language?
Tick the two correct answers.

\textbf{Correct Answer(s):}
\begin{itemize}
    \item MIPS assembly instructions typically use three operands for R-type instructions.
    \item The register \$0 always holds the value zero in MIPS.
\end{itemize}

\textbf{Incorrect Answers:}
\begin{itemize}
    \item MIPS uses a single memory space for both instructions and data.
    \item In MIPS, memory access is permitted in any instruction type.
    \item MIPS does not support pipelined execution.
\end{itemize}

\textbf{Solution \& Feedback:}\\
MIPS uses a single memory space for both instructions and data.\\
This is incorrect. MIPS uses a modified Harvard architecture with separate interfaces for instructions and data, allowing parallel access.\\

In MIPS, memory access is permitted in any instruction type.\\
This is incorrect because memory access in MIPS is restricted to load (\texttt{lw}) and store (\texttt{sw}) instructions only, aligning with its load/store architecture principle.\\

MIPS does not support pipelined execution.\\
This is incorrect. MIPS actively uses pipelining to enhance execution speed by overlapping instruction stages.

\subsection{Variation Question 3.2}

Which of the following statements correctly describe aspects of the MIPS architecture and assembly language?

Tick the two correct answers.

\textbf{Correct Answer(s):}
\begin{itemize}
    \item MIPS uses two separate memory spaces for instructions and data.
    \item MIPS assembly instructions typically use three operands for R-type instructions.
\end{itemize}

\textbf{Incorrect Answers:}
\begin{itemize}
    \item In MIPS, memory access is permitted in any instruction type.
    \item The register \$s0 always holds the value zero in MIPS.
    \item MIPS does not support pipelined execution.
\end{itemize}

\textbf{Solution \& Feedback:}\\
In MIPS, memory access is permitted in any instruction type.\\
This is incorrect because memory access in MIPS is restricted to load (\texttt{lw}) and store (\texttt{sw}) instructions only, aligning with its load/store architecture principle.\\

The register \$s0 always holds the value zero in MIPS.\\
This is incorrect. The \$s0 register is used to store intermediate values during computations.\\

MIPS does not support pipelined execution.\\
This is incorrect. MIPS actively uses pipelining to enhance execution speed by overlapping instruction stages.

\subsection{Variation Question 3.3}

Which of the following statements correctly describe aspects of the MIPS architecture and assembly language?

Tick the two correct answers.

\textbf{Correct Answer(s):}
\begin{itemize}
    \item MIPS uses two separate memory spaces for instructions and data.
    \item Memory access in MIPS is restricted to load (\texttt{lw}) and store (\texttt{sw}) instructions only, aligning with its load/store architecture principle.
\end{itemize}

\textbf{Incorrect Answers:}
\begin{itemize}
    \item MIPS assembly instructions typically use three operands for J-type instructions.
    \item The register \$s0 always holds the value zero in MIPS.
    \item MIPS does not support pipelined execution.
\end{itemize}

\textbf{Solution \& Feedback:}\\
MIPS assembly instructions typically use three operands for J-type instructions.\\
This is incorrect. J-type only requires Opcode and Value.\\

The register \$s0 always holds the value zero in MIPS.\\
\textcolor{red}{
This is incorrect. The \$s0 register is used to store intermediate values during computations.\\
}

MIPS does not support pipelined execution.\\
This is incorrect. MIPS actively uses pipelining to enhance execution speed by overlapping instruction stages.

\subsection{Variation Question 4.1}

Which of the following statements about compilation and interpretation is false?

\textbf{Correct Answer(s):}
\begin{itemize}
    \item In general, compilation is less computationally expensive than interpretation.
\end{itemize}

\textbf{Incorrect Answers:}
\begin{itemize}
    \item In general, executing an interpreted program is slower than executing a compiled program.
    \item Virtual machines involve both compilation and interpretation.
    \item In general, interpretation facilitates debugging better than compilation.
\end{itemize}

\textbf{Solution \& Feedback:}\\
It is true that, in general, executing an interpreted program is slower than executing a compiled program.\\
It is true that virtual machines involve both compilation and interpretation.\\
It is true that, generally, interpretation facilitates debugging better than compilation.\\
It is not true that compilation is generally less computationally expensive than interpretation – compilation is generally more expensive than interpretation.

\subsection{Variation Question 4.2}

Which of the following statements about compilation and interpretation is true?

\textbf{Correct Answer(s):}
\begin{itemize}
    \item In general, executing an interpreted program is slower than executing a compiled program.
\end{itemize}

\textbf{Incorrect Answers:}
\begin{itemize}
    \item In general, compilation is less computationally expensive than interpretation.
    \item In general, compilation facilitates debugging better than interpretation.
    \item In general, Java Virtual Machine does not include any type of compilation and performs interpretation only.
\end{itemize}

\textbf{Solution \& Feedback:}\\
It is not true that compilation is generally less computationally expensive than interpretation – compilation is generally more expensive than interpretation.\\
It is not true that compilation facilitates debugging better than interpretation – interpretation facilitates debugging better than compilation.\\
It is not true that Java Virtual Machine does not include any type of compilation and performs interpretation only. JVM invokes the Just-in-Time compilation process, whenever a code segment is executed repeatedly.\\
It is true that executing an interpreted program is generally slower than executing a compiled program.

\subsection{Variation Question 4.3}

Which of the following statements about compilation and interpretation is false?

\textbf{Correct Answer(s):}
\begin{itemize}
    \item In general, Java Virtual Machine does not include any type of compilation and performs interpretation only.
\end{itemize}

\textbf{Incorrect Answers:}
\begin{itemize}
    \item In general, executing an interpreted program is slower than executing a compiled program.
    \item In general, interpretation facilitates debugging better than compilation.
    \item In general, compilation is more computationally expensive than interpretation.
\end{itemize}

\textbf{Solution \& Feedback:}\\
It is true that, in general, executing an interpreted program is slower than executing a compiled program.\\
It is true that, in general, compilation is more computationally expensive than interpretation.\\
It is true that executing an interpreted program is generally slower than executing a compiled program.\\
It is not true that Java Virtual Machine does not include any type of compilation and performs interpretation only. JVM invokes the Just-in-Time compilation process, whenever a code segment is executed repeatedly.

% --- End of Part 2: Next comes IEEE 754 (Variation 5.x) onwards
\subsection{Variation Question 5.1}

Given the following number in IEEE 754 (32-bit) floating point format with 8-bit exponent, calculate the decimal real number it represents.

$1\ 1000\ 0111\ 001\ 0001\ 0001\ 0000\ 0000\ 0000$

\textbf{Correct Answer(s):}
\begin{itemize}
    \item $-290.125$
\end{itemize}

\textbf{Incorrect Answers:}
\begin{itemize}
    \item $-145.0625$ (on subtracting 128 from exponent)
    \item $-72.015625$ (2’s complement exponent + 127)
    \item $-0.008853912$ (sign-and-magnitude exponent)
\end{itemize}

\textbf{Solution \& Feedback:}\\
Sign bit: 1 (-ve)\\
Biased Exponent: (1000 0111)$_2$ = 135\\
Real Exponent: 135 – 127 = 8\\
Stored Mantissa: (001 0001 0001 0000 0000 0000)$_2$\\
Actual Mantissa: (1.00100010001)$_2$\\
Normalized Representation: $1.00100010001 \times 2^8$\\
Binary Number: 100100010.001\\
Decimal Value: $-290.125$

\subsection{Variation Question 5.2}

Given the following number in IEEE 754 (32-bit) floating point format with 8-bit exponent, calculate the decimal real number it represents.

$1\ 1000\ 0100\ 101\ 0000\ 1000\ 1000\ 0000\ 0000$

\textbf{Correct Answer(s):}
\begin{itemize}
    \item $-52.1328125$
\end{itemize}

\textbf{Incorrect Answers:}
\begin{itemize}
    \item $-26.06640625$ (on subtracting 128 from exponent)
    \item $-13.033203125$ (2’s complement exponent + 127)
    \item $-0.101821899$ (sign-and-magnitude exponent)
\end{itemize}

\textbf{Solution \& Feedback:}\\
Sign bit: 1 (-ve)\\
Biased Exponent: (1000 0100)$_2$ = 132\\
Real Exponent: 132 – 127 = 5\\
Stored Mantissa: (101 0000 1000 1000 0000 0000)$_2$\\
Actual Mantissa: (1.101000010001)$_2$\\
Normalized Representation: $1.101000010001 \times 2^5$\\
Binary Number: 110100.0010001\\
Decimal Value: $-52.1328125$

\subsection{Variation Question 5.3}

Given the following number in IEEE 754 (32-bit) floating point format with 8-bit exponent, calculate the decimal real number it represents.

$1\ 1000\ 1001\ 100\ 0101\ 1100\ 0110\ 1000\ 0000$

\textbf{Correct Answer(s):}
\begin{itemize}
    \item $-1582.203125$
\end{itemize}

\textbf{Incorrect Answers:}
\begin{itemize}
    \item $-791.1015625$ (on subtracting 128 from exponent)
    \item $-395.55078125$ (2’s complement exponent + 127)
    \item $-0.003017813$ (sign-and-magnitude exponent)
\end{itemize}

\textbf{Solution \& Feedback:}\\
Sign bit: 1 (-ve)\\
Biased Exponent: (1000 1001)$_2$ = 137\\
Real Exponent: 137 – 127 = 10\\
Stored Mantissa: (100 0101 1100 0110 1000 0000)$_2$\\
Actual Mantissa: (1.1000101110001101)$_2$\\
Normalized Representation: $1.1000101110001101 \times 2^{10}$\\
Binary Number: 11000101110.001101\\
Decimal Value: $-1582.203125$

% ---------------------------------------

\subsection{Variation Question 6.1}

Initial Memory Values (decimal):\\
Memory address 1000: 5\\
Memory address 2000: 10\\
Memory address 3000: 15\\
Memory address 4000: 20

Instructions:
\begin{itemize}
    \item ADD 1000, 2000, 3000
    \item ADD 4000, 1000, 2000
    \item MUL 1000, 4000, 3000
\end{itemize}

\textbf{Correct Answer(s):}
\begin{itemize}
    \item 525
\end{itemize}

\textbf{Incorrect Answers:}
\begin{itemize}
    \item 5
    \item 25
    \item 35
\end{itemize}

\textbf{Solution \& Feedback:}\\
ADD 1000, 2000, 3000 $\Rightarrow$ Address1000 = 10+15 = 25\\
ADD 4000, 1000, 2000 $\Rightarrow$ Address4000 = 25+10 = 35\\
MUL 1000, 4000, 3000 $\Rightarrow$ Address1000 = 35*15 = 525

% --- continuing with 6.2, 6.3, 7.x, 8.x, etc., next immediately
\subsection{Variation Question 6.2}

Initial Memory Values (decimal):\\
Memory address 1000: 0\\
Memory address 2000: 15\\
Memory address 3000: 2\\
Memory address 4000: 10

Instructions:
\begin{itemize}
    \item ADD 1000, 2000, 3000
    \item ADD 4000, 1000, 2000
    \item MUL 1000, 4000, 3000
\end{itemize}

\textbf{Correct Answer(s):}
\begin{itemize}
    \item 64
\end{itemize}

\textbf{Incorrect Answers:}
\begin{itemize}
    \item 0
    \item 32 
    \item 70 
\end{itemize}

\textbf{Solution \& Feedback:}\\
ADD 1000, 1000, 2000 $\Rightarrow$ Address1000 = 15+2 = 17\\
ADD 4000, 1000, 3000 $\Rightarrow$ Address4000 = 17+15 = 32\\
MUL 1000, 4000, 3000 $\Rightarrow$ Address1000 = 32*2 = 64\\
(Note: It seems that the final correct answer is given as 64 in your quiz. Verify memory changes and multiplication steps if needed.)

\subsection{Variation Question 6.3}

Initial Memory Values (decimal):\\
Memory address 1000: 15\\
Memory address 2000: 5\\
Memory address 3000: 20\\
Memory address 4000: 10

Instructions:
\begin{itemize}
    \item ADD 1000, 2000, 3000
    \item ADD 4000, 1000, 2000
    \item MUL 1000, 4000, 3000
\end{itemize}

\textbf{Correct Answer(s):}
\begin{itemize}
    \item 600
\end{itemize}

\textbf{Incorrect Answers:}
\begin{itemize}
    \item 15
    \item 30
    \item 700
\end{itemize}

\textbf{Solution \& Feedback:}\\
ADD 1000, 1000, 2000 $\Rightarrow$ Address1000 = 5+20 = 25\\
ADD 4000, 2000, 3000 $\Rightarrow$ Address4000 = 25+5 = 35\\
ADD 3000, 4000, 3000 $\Rightarrow$ Address3000 = 30$\times$20 = 40

(Note: Cross-check logic based on actual addition in registers.)

\subsection{Variation Question 7.1}
Consider the following MIPS code:

\begin{lstlisting}
    li \$t0, 1
    li \$t1, &x
Loop: 
    beq \$t2, 0, exit
    mul $t0, $t0, $t1
    sub \$t2, \$t2, 1
    j loop
exit:
    sw $t0, &z
\end{lstlisting}
Assuming initial values of x = 3 and y = 4, what will be the
final values of z when the code completes its execution.
What value is left in \$t2 after the above MIPS code executes?
Tick the correct answer.
Answer choices:

\textbf{Solution and Feedback:}
In this question, the initial value of \$t0 is 3, the initial value of \$t1 is 4, and the initial value of \$t2 is 1. 
In the loop:
- Multiply \$t2 by \$t0
- Subtract 1 from \$t1
- Continue looping until \$t1 is 0

Initially, \$t2 = 1.

First loop:
\$t2 = 1 * 3 = 3
\$t1 = 3

Second loop:
\$t2 = 3 * 3 = 9
\$t1 = 2

Third loop:
\$t2 = 9 * 3 = 27
\$t1 = 1

Fourth loop:
\$t2 = 27 * 3 = 81
\$t1 = 0

Thus, the final value in \$t2 is 81.

\subsection{Variation Question 7.2}

\textbf{Question} \\
Consider the following MIPS assembly code:
\begin{lstlisting}
    li $t0, 1
    lw $t1, &x
    lw $t2, &y
loop:
    beq $t2, 0, exit
    mul $t0, $t0, $t1
    sub $t2, $t2, 1
    j loop
exit:
    sw $t0, &z
\end{lstlisting}

Assuming initial values of x = -4 and y = 3, what will be the final values of z when the code completes its execution.\\

\textbf{Correct Answer(s)} : -64  \\
\noindent \textbf{Solution \& Feedback:} \\
\noindent The MIPS Assembly code computes the exponentiation operation, where x is the base and y contains the exponent. The result is stored in z.\\
\noindent The following tables shows the values of registers during execution:

\begin{tabular}{|c|c|c|}
\hline
\textbf{\$t0 (z)} & \textbf{\$t1 (x)} & \textbf{\$t2 (y)} \\
\hline
1 & -4 & 3 \\
\hline
-4 & -4 & 2 \\
\hline
16 & -4 & 1 \\
\hline
-64 & -4 & 0 \\
\hline
\end{tabular}

\noindent Once \$t2 reaches 0, the loop terminates, and the result (\textbf{-64}) is stored in z.

\subsection{Variation Question 7.3}
\textbf{Variation Question} \\
Consider the following MIPS assembly code: \\
\begin{lstlisting}
    li $t0, 1
    lw $t1, &x
    lw $t2, &y
loop:
    beq $t2, 0, exit
    mul $t0, $t0, $t1
    sub $t2, $t2, 1
    j loop
exit:
    sw $t0, &z
\end{lstlisting}
Assuming initial values of x = -2 and y = 5, what will be the final values of z when the code completes its execution.\\


\textbf{Correct Answer(s)} : -32  \\

\noindent \textbf{Solution \& Feedback:}
\noindent The MIPS Assembly code computes the exponentiation operation, where x is the base and y contains the exponent. The result is stored in z.
\noindent The following tables shows the values of registers during execution:

\begin{tabular}{|c|c|c|}
\hline
\textbf{\$t0 (z)} & \textbf{\$t1 (x)} & \textbf{\$t2 (y)} \\
\hline
1 & -2 & 5 \\
\hline
-2 & -2 & 4 \\
\hline
4 & -2 & 3 \\
\hline
-8 & -2 & 2 \\
\hline
16 & -2 & 1 \\
\hline
-32 & -2 & 0 \\
\hline
\end{tabular}

\noindent Once \$t2 reaches 0, the loop terminates, and the result (\textbf{-32}) is stored in z.

\subsection{Variation Question 8.1}
\textbf{Question} \\
Evaluate the following RPN expressions using a stack and indicate the correct answer:\\
3 5 8 * 7 + 4 - +
\textbf{Note:} All the values in the expression above are single digit numbers.\\

\textbf{Correct Answer(s)} : 46 

\noindent \textbf{Solution \& Feedback:}\\
\noindent We can evaluate the RPN expression using the following steps:\\

\begin{tabular}{|c|c|c|}
\hline
\textbf{Step \#} & \textbf{Input} & \textbf{Stack (Top on Right)} \\
\hline
1 & 3 & 3 \\
\hline
2 & 5 & 3 5 \\
\hline
3 & 8 & 3 5 8 \\
\hline
4 & * & 3 40 \\
\hline
5 & 7 & 3 40 7 \\
\hline
6 & + & 3 47 \\
\hline
7 & 4 & 3 47 4 \\
\hline
8 & - & 3 43 \\
\hline
9 & + & 46 \\
\hline
\end{tabular}

\noindent The final answer is \textbf{46}.

\subsection{Variation Question 8.2}
\textbf{Question} \\
Evaluate the following RPN expressions using a stack and indicate the correct answer:\\
9 3 + 8 2 * 7 - + \\
\textbf{Note:} All the values in the expression above are single digit numbers. \\ 
\textbf{Correct Answer(s)} 21 \\
\textbf{Solution \& Feedback:} \\
We can evaluate the RPN expression using the following steps: \\
\begin{tabular}{|c|c|c|}
\hline
\textbf{Step \#} & \textbf{Input} & \textbf{Stack (Top on Right)} \\ \hline
1 & 9 & 9 \\ \hline
2 & 3 & 9 3 \\ \hline
3 & + & 12 \\ \hline
4 & 8 & 12 8 \\ \hline
5 & 2 & 12 8 2 \\ \hline
6 & * & 12 16 \\ \hline
7 & 7 & 12 16 7 \\ \hline
8 & - & 12 9 \\ \hline
9 & + & 21 \\ \hline
\end{tabular}

The final answer is \textbf{21}. \\
\textbf{Answer:} 21 \\

\subsection{Variation Question 8.3}
\textbf{Question} \\ 
Evaluate the following RPN expressions using a stack and indicate the correct answer:
8 4 6 - * 3 9 + + \\
\textbf{Note:} All the values in the expression above are single digit numbers. \\ 
\textbf{Correct Answer(s)}
-4 \\
\textbf{Solution \& Feedback:} \\ 
We can evaluate the RPN expression using the following steps: \\
\begin{tabular}{|c|c|c|}
\hline
\textbf{Step \#} & \textbf{Input} & \textbf{Stack (Top on Right)} \\ \hline
1 & 8 & 8 \\ \hline
2 & 4 & 8 4 \\ \hline
3 & 6 & 8 4 6 \\ \hline
4 & - & 8 -2 \\ \hline
5 & * & -16 \\ \hline
6 & 3 & -16 3 \\ \hline
7 & 9 & -16 3 9 \\ \hline
8 & + & -16 12 \\ \hline
9 & + & -4 \\ \hline
\end{tabular}

The final answer is \textbf{-4}.
\textbf{Answer:} $-4$

\subsection{Variation 9.1}
\textbf{Question} \\
Your friend tells you that they have come up with a new fixed-length representation of text characters using binary. They do not tell you anything about how the representation works.\\
They write the following secret message to you using this representation: \\
\texttt{11000110100110}
They then give you five options for what the message says in English. Which of the five options is represented by the secret message?

\textbf{Correct Answer(s):} PRESSES \\
\textbf{Incorrect Answers:} TATTOOS, MONSOON, PRECEED, RACCOON

\textbf{Solution \& Feedback:}
Each of the answers is 7 characters long, and the word is 14-bits long: therefore each character is 2-bits long. As each character is two bits long, we can only represent $2^2 = 4$ different characters using this representation: that eliminates PRECEED and RACCOON.
The candidates are now TATTOOS, PRESSES or MONSOON. Repeated 2-bit codes can be used to identify which is the correct word:
\texttt{11 00 01 10 10 01 10}
'10' is repeated in the 4\textsuperscript{th}, 5\textsuperscript{th} and 7\textsuperscript{th} characters; therefore, the word must be PRESSES.

\subsection{Variation 9.2}
\textbf{Question} \\
Your friend tells you that they have come up with a new fixed-length representation of text characters using binary. They do not tell you anything about how the representation works.
They write the following secret message to you using this representation:\\
\texttt{1100011010}\\
They then give you five options for what the message says in English. Which of the five options is represented by the secret message?

\textbf{Correct Answer(s):} SMALL \\
\textbf{Incorrect Answers:} GOING, RODEO, SIDED, SNARL

\textbf{Solution \& Feedback:}\\
Each of the answers is 5 characters long, and the word is 10-bits long: therefore each character is 2-bits long. As each character is two bits long, we can only represent $2^2 = 4$ different characters using this representation: that eliminates SNARL.\\
The candidates are now GOING, RODEO, SMALL or SIDED. Repeated 2-bit codes can be used to identify which is the correct word:\\
\texttt{11 00 01 10 10}\\
'10' is doubled in the 4\textsuperscript{th} and 5\textsuperscript{th} characters; therefore, the word must be SMALL.

\subsection{Variation 9.3}
\textbf{Question} \\
Your friend tells you that they have come up with a new fixed-length representation of text characters using binary. They do not tell you anything about how the representation works.
They write the following secret message to you using this representation: 
\texttt{110001101000}\\
They then give you five options for what the message says in English. Which of the five options is represented by the secret message?

\textbf{Correct Answer(s):} STREET \\
\textbf{Incorrect Answers:} ACCESS, COMMON, SNAPPY, CHATTY

\textbf{Solution \& Feedback:}
Each of the answers is 6 characters long, and the word is 12-bits long: therefore each character is 2-bits long. As each character is two bits long, we can only represent $2^2 = 4$ different characters using this representation: that eliminates SNAPPY and CHATTY.\\
The candidates are now ACCESS, COMMON or STREET. Note that the double letters in both can help identify which is the correct word:\\
\texttt{11 00 01 10 10 00}\\
'10' is doubled in the 4\textsuperscript{th} and 5\textsuperscript{th} characters; therefore, the word must be STREET.\\

\subsection{Variation Question 10.1}
Consider the following MIPS assembly code:
\begin{lstlisting}
    lw $t1, &a
    lw $t2, &b
    mul $t3, $t1, $t2
L1:
    bne $t1, $t2, L2
    j exit
L2:
    bgt $t1, $t2, L3
    sub $t2, $t2, $t1
    j L1
L3:
    sub $t1, $t1, $t2
    j L1
exit:
    div $t3, $t3, $t1
    sw $t3, &c
\end{lstlisting}
Assuming initial values of a = 15 and b = 25, what will be the value of c when the code completes its execution? \\
\textbf{Correct Answer(s)} \\
75 \\

\textbf{Solution \& Feedback:} \\
The above code implements the LCM algorithms using the following: \\
\(LCM (a, b) = (a \times b) / GCD (a, b) \)
Initially, the code computes the GCD using the following algorithm:

\begin{lstlisting}
int gcd(int a, int b)
{
    while (a != b)
        if (a > b)
            a = a - b;
        else
            b = b - a;
    return a;
}
\end{lstlisting}

For a = 15 and b = 25, the following table shows how the register values will be updated:

\begin{center}
\begin{tabular}{|c|c|c|}
\hline
$t1$ (a) & $t2$ (b) & $t3$ (c) \\
\hline
15 & 25 & 375 \\
15 & 10 & 375 \\
5  & 10 & 375 \\
5  & 5  & 375 \\
\hline
\end{tabular}
\end{center}

When the values of registers $t1$ and $t2$ become equal, the code jumps to exit and computes the following:

\[
\text{LCM} (a, b) = (a * b) / \text{GCD} (a, b) = 375 / 5 = \mathbf{75}
\]
LCM (Least Common Multiple) is the smallest positive integer that is divisible by all the given integers, while GCD (Greatest Common Divisor) is the largest positive integer that divides all the given integers without a remainder

\subsection{Variation Question 10.2}
Consider the following MIPS assembly code:\\
\begin{lstlisting}
    lw $t1, &a
    lw $t2, &b
    mul $t3, $t1, $t2
L1:
    bne $t1, $t2, L2
    j exit
L2:
    bgt $t1, $t2, L3
    sub $t2, $t2, $t1
    j L1
L3:
    sub $t1, $t1, $t2
    j L1
exit:
    div $t3, $t3, $t1
    sw $t3, &c
\end{lstlisting}
Assuming initial values of a = 16 and b = 24, what will be the value of c when the code completes its execution? \\
\textbf{Correct Answer(s)} \\
48 \\
\textbf{Solution \& Feedback:} \\
The above code implements the LCM algorithms using the following: \\
\(LCM (a, b) = (a * b) / GCD (a, b) \)
Initially, the code computes the GCD using the following algorithm:
\begin{lstlisting}
int gcd(int a, int b)
{
    while (a != b)
        if (a > b)
            a = a - b;
        else
            b = b - a;
    return a;
}
\end{lstlisting}
For a = 15 and b = 25, the following table shows how the register values will be updated:

\begin{center}
\begin{tabular}{|c|c|c|}
\hline
$t1$ (a) & $t2$ (b) & $t3$ (c) \\
\hline
16 & 24 & 384 \\
16 & 8  & 384 \\
8  & 8  & 384 \\
\hline
\end{tabular}
\end{center}
When the values of registers $t1$ and $t2$ become equal, the code jumps to exit and computes the following:
\[
\text{LCM} (a, b) = (a * b) / \text{GCD} (a, b) = 384 / 8 = \mathbf{48}
\]
\subsection{Variation 10.3}
\textbf{Question} 
Consider the following MIPS assembly code:
\begin{lstlisting}
    lw $t1, &a
    lw $t2, &b
    mul $t3, $t1, $t2
L1:
    bne $t1, $t2, L2
    j exit
L2:
    bgt $t1, $t2, L3
    sub $t2, $t2, $t1
    j L1
L3:
    sub $t1, $t1, $t2
    j L1
exit:
    div $t3, $t3, $t1
    sw $t3, &c
\end{lstlisting}
Assuming initial values of $a=9$ and $b=15$, what will be the value of $c$ when the code completes its execution?\\
\textbf{Correct Answer(s)}: 45 \\
\textbf{Incorrect Answers}: 48, 75, 135, 375, 384
\noindent\textbf{Solution \& Feedback:} \\
The above code implements the LCM algorithms using the following: \\
LCM $(a,b) = (a*b) / \text{GCD}(a,b)$
Initially, the code computes the GCD using the following algorithm:
\begin{lstlisting}
int gcd(int a, int b)
{
    while (a != b)
        if (a > b)
            a = a - b;
        else
            b = b - a;
    return a;
}
\end{lstlisting}
For $a=15$ and $b=25$, the following table shows how the register values will be updated:
\begin{center}
\begin{tabular}{|c|c|c|}
\hline
$t1$ (a) & $t2$ (b) & $t3$ (c) \\
\hline
9 & 15 & 135 \\
9 & 6  & 135 \\
3 & 6  & 135 \\
3 & 3  & 135 \\
\hline
\end{tabular}
\end{center}
When the values of registers $t1$ and $t2$ become equal, the code jumps to exit and computes the following:

\[
\text{LCM}(a,b) = (a*b)/\text{GCD}(a,b) = 135/3 = \mathbf{45}
\]
